\worksheettitle
  {M02-R. Показания концов}
  {\modulemark{M02-R} Коррекционная карточка}

\studentnameblock

\begin{warningbox}[Назовите ошибку]
  Показание правого конца совпадает с длиной только тогда, когда левый конец
  находится у нуля. В остальных случаях нужно найти разность показаний.
\end{warningbox}

\begin{practicebox}[Шаг 1. Начало у нуля]
  Левый конец у $0$, правый у $4\text{ см }7\text{ мм}$.
  Длина равна \answerblank[45mm]. Почему вычитать ничего не нужно?
  \answerblank[75mm]
\end{practicebox}

\begin{practicebox}[Шаг 2. Начало сдвинуто]
  Левый конец у $2\text{ см}$, правый у $6\text{ см }5\text{ мм}$.
  \[
    \answerblank[34mm]-\answerblank[34mm]=\answerblank[40mm].
  \]
\end{practicebox}

\begin{practicebox}[Шаг 3. Исправьте]
  Ученик записал: «От $3\text{ см }1\text{ мм}$ до
  $8\text{ см }9\text{ мм}$ длина равна $8\text{ см }9\text{ мм}$».

  Причина: \answerblank[95mm]

  Исправление: \answerblank[95mm]
\end{practicebox}

\begin{checkpointbox}[Повторная проверка]
  Концы нового отрезка находятся у $1\text{ см }6\text{ мм}$ и
  $7\text{ см }3\text{ мм}$. Найдите длину и объясните действие.
  \writinglines{2}
\end{checkpointbox}
